\section{Open Questions}

Five questions the replication and re-reading of the paper leave genuinely open. These are
not rhetorical --- each has a concrete next step and would materially change the interpretation
of the paper if answered.

\subsection*{OQ1: Is the sub-minute temporal resolution real or asserted?}
\textbf{Basis.} The paper's central mechanistic claim --- that early NBS1 recruitment is
dominated by an ``inner-focus'' population binding directly to DSB ends \emph{before} the
H2AX/MDC1 scaffold assembles --- depends on cleanly resolving the first
$\sim 30\text{--}60$\,s post-irradiation. But the live-cell beamline setup requires a
shutter-and-stage sequence between irradiation and confocal acquisition, and the paper
does not publish the actual first-frame latency, frame interval, or per-frame photobleach
correction. The ODE model can be fit to whatever early-time signal the imaging returns,
so ``the model agrees'' is compatible with the imaging having $\pm 10\text{--}20$\,s
effective time resolution rather than the sub-minute precision the interpretation implies.

\textbf{Next steps.}
(1) Request from the GSI / TU Darmstadt authors the raw beamline log with per-frame timestamps
for at least one recruitment curve at each of low, mid, and high LET.
(2) Re-fit the ODE model with the first data point held out and see whether the inner-focus
fraction at $t<60$\,s is still constrained or becomes a free parameter.
(3) If timestamps are unavailable, publish an explicit sensitivity analysis: how do the fitted
rate constants and the inner-vs-outer decomposition shift when first-frame latency is varied
from 5 to 60\,s?

\subsection*{OQ2: Does the cylindrical-track / homogeneous-nucleus geometry bias the fit?}
\textbf{Basis.} The model treats an ion track as a well-defined cylinder of DSBs and the
nucleus as a homogeneous cylinder of ATM, MRN, and H2AX substrate. In reality the track has
a radial dose profile dominated by $\delta$-electron spurs (falling as $\sim 1/r^2$ out to
$\mu$m radii), and chromatin is not spatially uniform. The fitted rate constants absorb any
spatial mismatch, so ``a good fit with a single global parameter set'' is compatible with the
true spatial DSB distribution being materially different from cylindrical.

\textbf{Next steps.}
(1) Convolve the assumed cylindrical track with a published radial-dose kernel
(Kiefer-Straaten or Chatterjee-Schaefer) and re-fit; compare rate constants to the paper's
published values and quote the shift.
(2) Run the model with a euchromatin/heterochromatin bimodal substrate concentration and
check whether inner-focus fraction predictions are robust.
(3) Compare model predictions against independent super-resolution imaging (e.g.\ STORM of
$\gamma$H2AX or MDC1) to test the cylindrical assumption directly.

\subsection*{OQ3: Are the 7 optimized rate constants uniquely determined, or degenerate?}
\textbf{Basis.} The paper reports point-value rate constants with no confidence intervals,
no bootstrap, no profile-likelihood, and no sensitivity analysis. Nine reactions, ten rate
constants, 4 initial concentrations, 12 per-panel scaling factors, and 16 recruitment
curves --- the fit is likely underdetermined in several directions (paired forward/reverse
pairs $k_{5f}/k_{5r}$ and $k_{6f}/k_{6r}$ are prime suspects). Without raw data we cannot
check this. If the fit is not unique, any mechanistic interpretation of individual rate
values is not evidence-supported.

\textbf{Next steps.}
(1) Ask the authors for raw recruitment CSVs and re-run Nelder-Mead + bootstrap; report
95\% CI on each rate constant.
(2) If raw data is unavailable, run a profile-likelihood sweep against the digitized figure
points to bound identifiability locally.
(3) Publish a sensitivity table: for each of the 12 headline predictions, which rate
constants dominate the derivative?

\subsection*{OQ4: Does the mechanism generalize beyond heavy ions and beyond U2OS?}
\textbf{Basis.} The inner-vs-outer decomposition is calibrated on high-LET heavy ions,
where DSBs are spatially clustered along a well-defined path, in a single p53-mutant cancer
cell line (U2OS). At low LET (X-rays, protons $<10$\,keV/$\mu$m) DSBs are more spatially
isolated and the ``inner focus'' signal per unit fluence should be much smaller. The paper
touches this via X-ray controls but does not systematically vary radiation quality across
the LET-continuum minimum, and does not test other cell lines. Applicability at
proton-therapy-relevant LETs and in non-cancer cell lines is genuinely open.

\textbf{Next steps.}
(1) Run the published ODE model at LET = 1, 2, 5, 10\,keV/$\mu$m (proton clinical range)
and predict inner-focus fraction; compare against any existing proton-beam FRAP data
(Antonelli et al., Chaudhary et al.).
(2) Solicit collaboration with a proton-therapy facility to run the same NBS1-GFP FRAP
protocol at clinically-relevant LETs.
(3) Re-derive the rate constants in a non-transformed cell line (RPE-1 or primary human
fibroblasts) and check whether inner-focus dominance at high LET holds.

\subsection*{OQ5: How does the network couple to chromatin state?}
\textbf{Basis.} The model treats H2AX, MRN, and ATM as spatially uniform substrate pools,
but the DDR is known to be chromatin-context-dependent: repair kinetics differ between
heterochromatin and euchromatin (Goodarzi \& Jeggo), ATM is required specifically for
heterochromatic DSB repair, and 53BP1/BRCA1 pathway choice depends on cell-cycle phase.
Because CK2 phosphorylation of MDC1 is itself chromatin-context-modulated, the
inner-vs-outer split may be systematically biased in either direction depending on the
chromatin fraction traversed by the ion track.

\textbf{Next steps.}
(1) Extend the ODE model to a two-compartment (heterochromatin + euchromatin) substrate with
distinct rate constants; fit against the same curves; test AIC/BIC preference.
(2) Use Hi-C / ChIP-seq maps of U2OS to compute expected heterochromatin fraction along a
random ion track and predict panel-to-panel variance currently attributed to noise.
(3) Run imaging with an HP1$\alpha$-mCherry heterochromatin marker to directly measure
whether inner-focus NBS1 kinetics differ inside vs.\ outside heterochromatin.
